% =============================================================
%  SECTION 1 — INTRODUCTION  (revised)
%  Target: ~1 page / ~500 words / 5 focused paragraphs
%  All marketing language removed; contributions are defensible
% =============================================================
\section{Introduction}

% --- P1: Real-world problem (concrete, specific) ---
Enterprise procurement workflows in emerging markets generate document inputs
that are predominantly unstructured: handwritten order slips, photographed
receipts, and verbal purchase requests transmitted over phone or messaging
applications. In the Vietnamese retail and distribution sector, this pattern is
especially pronounced. A significant share of small and medium enterprises
(SMEs) record transactions on paper invoices and communicate orders verbally,
often mixing Vietnamese with English product terminology and regional
abbreviations~\cite{chuong2019robotic}. Existing Robotic Process Automation
(RPA) deployments handle these inputs through manual exception queues, which
erode the efficiency gains that automation is intended to deliver.

% --- P2: Technical challenge (precise, no hyperbole) ---
Processing such inputs end-to-end with large language models (LLMs) is
technically feasible but operationally impractical at scale. Invoking a
multimodal LLM for every incoming document incurs per-token API costs and
round-trip latency that, across the high transaction volumes typical of
distribution centres, accumulate to prohibitive expense~\cite{reid2024gemini}.
At the same time, lightweight rule-based pipelines fail on the lexical
variability introduced by tonal Vietnamese, bilingual product names, and
inconsistent abbreviation conventions (e.g., \textit{cty tnhh} for
\textit{cong ty trach nhiem huu han})~\cite{le2024phowhisper}. Neither
extreme satisfies the accuracy and throughput requirements of production RPA.

% --- P3: Gap in existing literature ---
Prior work on document information extraction has focused on English-language,
high-resource settings~\cite{xu2020layoutlm} or on isolated components—ASR
for Vietnamese~\cite{le2024phowhisper, zhuo2025vietasr}, entity extraction
via LLM prompting~\cite{reid2024gemini}, and hybrid retrieval for product
search~\cite{reimers2019sentence, cormack2009reciprocal}—without
integrating these into a unified pipeline designed for enterprise-grade
deployment constraints. No publicly available system addresses the combination
of Vietnamese multimodal input, zero-shot entity extraction, and
catalog-scale entity resolution under latency and cost budgets relevant to
production RPA.

% --- P4: Our solution (what, not ``novel'' or ``innovative'') ---
This paper describes a modular, two-stage pipeline that separates
\emph{signal transcription} from \emph{semantic extraction}. In Stage~1,
modality-specific modules—a cloud OCR service for document images and
PhoWhisper~\cite{le2024phowhisper} for spoken Vietnamese—convert raw inputs
into a shared textual representation. In Stage~2, a single LLM performs
zero-shot structured extraction of merchants, products, quantities, and units
from the transcribed text. Extracted entities are then resolved against an
enterprise catalog through a heterogeneous resolution layer: fuzzy string
matching for merchant names and hybrid dense–sparse retrieval (Sentence-BERT +
BM25, fused via Reciprocal Rank Fusion) for product SKUs, with an optional
taxonomy-guided filter and a human-in-the-loop fallback for low-confidence
cases.

% --- P5: Contributions (defensible, scoped to what is actually done) ---
The contributions of this paper are:
\begin{enumerate}
    \item \textbf{System design}: A modular, eight-layer pipeline that
    processes Vietnamese visual and acoustic business documents into
    schema-compliant database records, with independent optimization of
    the extraction and resolution components.

    \item \textbf{Heterogeneous entity resolution}: A domain-driven
    resolution strategy that applies fuzzy lexical matching to merchant
    names and hybrid dense–sparse retrieval to product descriptions,
    motivated by the distinct error characteristics of each entity type.

    \item \textbf{Empirical evaluation}: Measurements of extraction
    accuracy (entity-level precision, recall, and F$_1$), product
    retrieval performance (MRR@$K$, Hit@1), ASR word error rate, and
    end-to-end throughput on a corpus of Vietnamese enterprise documents
    and audio recordings, with ablations across retrieval fusion
    strategies.

    \item \textbf{Deployment analysis}: A cost and latency breakdown of
    the modular architecture compared to an end-to-end LLM baseline,
    informing practical deployment decisions for RPA integration.
\end{enumerate}